Near-Earth Asteroids (NEAs) are small rocky bodies orbiting the Sun in the vicinity of Earth's orbit. A science mission designed to study a NEA target is operating a spacecraft that will make its closest approach to the target at a distance of $\Delta_c=100\,\mathrm m$. One of the instruments aboard the spacecraft is a CCD camera called the Asteroid Framer (AF), which can acquire images in the V-band at an effective wavelength of $\lambda=550\,\mathrm{nm}$. After the mission launched, the operations team conducted initial tests by acquiring images of the Earth using the AF camera with an aperture diameter of $D=10\,\mathrm{mm}$. At that time, the spacecraft was at a distance of $d=1.5\times10^6\,\mathrm{km}$ from Earth. The Earth was observed on the images within a photometric aperture of $\alpha_\oplus=60$ pixels in diameter.
\begin{description}
\item[(a)] Compute the pixel scale in arcseconds per pixel.
\item[(b)] Estimate the linear diameter of the smallest morphological feature that can be observed on the surface of the asteroid at the closest approach.
\end{description}
The saturation test on initial images showed Earth was just over-exposed at 10\,s exposure time and just under-exposed at 0.5\,s (under-exposed means signal-to-noise ratio $<5$). Using this data, the operations team can determine the exposure time interval allowed to observe the NEA target, which has the expected albedo of $A=0.10$. Ignore phase angle.

Assume that the detected signal is proportional to the object's surface brightness and the exposure time, and that the surface brightness is proportional to the albedo divided by the square of its distance from the Sun.
\begin{description}
\item[(c)] Determine the corresponding exposure time interval for which the planned observations of the NEA target at the closest approach will be successful without any under-exposure or over-exposure. At that time, the target will be located at $r_{\mathrm{NEA}}=1.20\,\mathrm{au}$ from the Sun.
\end{description}
